% Corte mecánico íntegro: parte_iv.tex:1--54.
% Residencia material del ensamblaje sucesor de 102 capítulos.
% Residencia de realización: capítulo 67.
\part{Geometría diferencial y realizaciones dimensionales}

\chapter{Curvatura y torsión de las curvas publicadas}

\section{Curva plana polar}

Sea
\[
 \gamma(\theta)=r(\theta)(\cos\theta,\sin\theta).
\]
Con primas respecto de \(\theta\), la velocidad y la curvatura son
\begin{align}
 v(\theta)&=\sqrt{r^2+(r')^2},\\
 \kappa_{\mathrm{F}}(\theta)&=
 \frac{r^2+2(r')^2-rr''}
 {\bigl(r^2+(r')^2\bigr)^{3/2}}.
 \label{espiral:eq:curvatura-polar}
\end{align}
La longitud de arco entre \(\theta_0\) y \(\theta_1\) es
\[
 s=\int_{\theta_0}^{\theta_1}
 \sqrt{r^2+(r')^2}\,\dd\theta.
\]
Estas magnitudes pertenecen a la curva publicada. No son el régimen TRIT ni
el cociclo radial.

\section{Especializaciones}

Para la espiral logarítmica \(r=r_0\e^{b\theta}\),
\[
 r'=br,\qquad r''=b^2r,
\]
y
\begin{equation}
 \kappa_{\mathrm{F}}(\theta)
 =\frac{1}{r(\theta)\sqrt{1+b^2}}.
 \label{espiral:eq:curvatura-log}
\end{equation}
El ángulo \(\psi\) entre el vector tangente y la dirección radial satisface
\[
 \tan\psi=\frac{r}{r'}=\frac1b,
\]
por lo que es constante. Ésta es la propiedad equiangular de la espiral
logarítmica.

\enlargethispage{4\baselineskip}
Para la espiral de Arquímedes \(r=a\theta\),
\[
 \kappa_{\mathrm{F}}(\theta)
 =\frac{\theta^2+2}{a(\theta^2+1)^{3/2}}.
\]
Para Fermat \(r=a\sqrt\theta\), la singularidad del origen exige declarar
\(\theta>0\). Para las ramas hiperbólica y lituus, el origen angular es una
asíntota y no un punto del dominio.
